%% kapa: exam=1 round=35 year=2024 period=1 form=A subject=real_estate subject_ko=부동산학원론 seq=36 number=116
%% source: Q-Net 감정평가사 기출문제 (cst00309) artlSeq=5213229 file=raw/1008/5213229/1교시 문제지 원본.pdf
%% conversion: kordoc 4.24.51 -> tools/gen1.py (auto); review=proofread-sonnet; flags=none
\begin{question}{116}
부동산투자분석에 관한 내용으로 옳지 않은 것은?
\begin{choices}
  \item 동일한 현금흐름을 가지는 투자안이라도 투자자의 요구수익률에 따라 순현재가치는 달라질 수 있다.
  \item 서로 다른 내부수익률을 가지는 두 자산에 동시에 투자하는 투자안의 내부수익률은 각 자산의 내부수익률을 더한 것과 같다.
  \item 동일한 투자안에 대해 내부수익률이 복수로 존재할 수 있다.
  \item 내부수익률법에서는 내부수익률과 요구수익률을 비교하여 투자의사결정을 한다.
  \item 투자규모에 차이가 나는 상호배타적인 투자안을 검토할 때, 순현재가치법과 수익성지수법을 통한 의사결정이 달라질 수 있다.
\end{choices}
\end{question}
\answer{2}
